\section*{Plan de cours : Introduction à GraphQL avec React Native (12 heures)}

\subsection*{Jour 1 : Fondamentaux de GraphQL (4h)}

\subsubsection*{Session 1.1 – Introduction & Installation (2h)}
\textbf{Objectif} : Comprendre les concepts de base et configurer l'environnement
\begin{itemize}
  \item Qu'est-ce que GraphQL ? (30 min)
  \item Comparaison REST vs GraphQL (30 min)
  \item Installation des outils (Node.js, npm, VS Code) (30 min)
  \item Initialisation d'un projet "Hello World" (30 min)
\end{itemize}

\vspace{0.3cm}
\subsubsection*{Session 1.2 – Schémas & Requêtes (2h)}
\textbf{Objectif} : Maîtriser les types et les requêtes simples
\begin{itemize}
  \item Types scalaires et objets (45 min)
  \item Écriture de requêtes et résolveurs (45 min)
  \item Exercice pratique : Calculatrice GraphQL (30 min)
\end{itemize}

\vspace{0.5cm}
\subsection*{Jour 2 : Projet Structuré & Mutations (4h)}

\subsubsection*{Session 2.1 – Structuration d’un projet GraphQL avec Prisma (2h)}
\textbf{Objectif} : Organiser un projet GraphQL de manière professionnelle
\begin{itemize}
  \item Restructuration du projet (dossiers resolvers, typedefs) (45 min)
  \item Configuration de Prisma et PostgreSQL (45 min)
  \item Seeders et premières requêtes (30 min)
\end{itemize}

\vspace{0.3cm}
\subsubsection*{Session 2.2 – Mutations & CRUD (2h)}
\textbf{Objectif} : Implémenter un CRUD complet
\begin{itemize}
  \item Mutations pour Contact (ajout / suppression) (1h)
  \item Tests avec Postman (30 min)
  \item Sécurisation des requêtes (30 min)
\end{itemize}

\vspace{0.5cm}
\subsection*{Jour 3 : Authentification & Consommation API (4h)}

\subsubsection*{Session 3.1 – Authentification JWT (2h)}
\textbf{Objectif} : Sécuriser l'API avec JWT
\begin{itemize}
  \item Théorie sur JWT et bcrypt (30 min)
  \item Mutations \texttt{register} et \texttt{login} (1h)
  \item Middleware d'authentification (30 min)
\end{itemize}

\vspace{0.3cm}
\subsubsection*{Session 3.2 – Consommation avec React Native (2h)}
\textbf{Objectif} : Consommer l’API depuis un client mobile
\begin{itemize}
  \item Configuration d'Apollo Client (45 min)
  \item Intégration des requêtes et mutations (1h)
  \item Exercice : Écran de login / dashboard (15 min)
\end{itemize}

\vspace{0.5cm}
\subsection*{Détails pédagogiques}
\begin{itemize}
  \item \textbf{Pratique} : 75\% de chaque session est dédiée à la mise en pratique.
  \item \textbf{Outils} : Postman, Prisma Studio, VS Code, Apollo Client.
  \item \textbf{Projet fil rouge} : Application "Gestion de Contacts".
\end{itemize}

\subsection*{Adaptabilité}
\begin{itemize}
  \item \textbf{Niveau débutant} : Accent mis sur les fondamentaux (Jour 1).
  \item \textbf{Niveau avancé} : Possibilité d’ajouter des notions avancées (caching, subscriptions, testing).
\end{itemize}

\subsection*{Conclusion}
Ce plan de formation garantit une montée en compétence progressive avec des résultats concrets dès le premier jour (API fonctionnelle, requêtes testables, début d’intégration mobile).

